\documentclass[11pt]{amsart}
\usepackage[T1]{fontenc}
\usepackage{lmodern}
\usepackage{parskip}
\usepackage[margin=1in]{geometry}
\usepackage{amsmath,amssymb,amsthm}
\usepackage[expansion=false]{microtype}
\usepackage[colorlinks=true,linkcolor=blue,citecolor=blue,urlcolor=blue]{hyperref}

\newtheorem{theorem}{Theorem}
\newtheorem{proposition}[theorem]{Proposition}
\newtheorem{lemma}[theorem]{Lemma}
\theoremstyle{remark}
\newtheorem{remark}[theorem]{Remark}
\newcommand{\eps}{\varepsilon}
\newcommand{\dt}{\,\frac{\mathrm{d}t}{t}}
\newcommand{\dx}{\,\frac{\mathrm{d}x}{x}}
\newcommand{\pos}[1]{\left(#1\right)^+}

\title[Daboussi's elementary proof, revisited]{Daboussi's elementary proof of the prime number theorem, revisited}
\author{Paul Pollack}
\address{Department of Mathematics, University of Georgia, Athens GA 30602}
\email{pollack@uga.edu}
\subjclass[2020]{Primary 11N05; Secondary 11N25}
\keywords{Prime number theorem, M\"obius function, smooth numbers, elementary methods}
\date{}

\begin{document}
\begin{abstract} In 1984, Daboussi gave an ``elementary'' proof of the prime number theorem whose central actors are the smooth numbers (numbers with no large prime factors). We explain a variant of his argument that might be considered even more elementary.
\end{abstract}
\maketitle

\section{Introduction}
Let $\mu$ denote the M\"obius function, and let $M(x):=\sum_{n\le x}\mu(n)$ be the corresponding summatory function. We will give a simple proof of the following 1897 theorem of von Mangoldt \cite{mangoldt}.
\begin{theorem}\label{thm:main}
$M(x)=o(x)$, as $x\to\infty$.
\end{theorem}

Around the turn of the 20th century, Landau showed that Theorem \ref{thm:main} is ``elementarily equivalent'' to the prime number theorem (PNT). (His 1899 doctoral thesis proves Theorem \ref{thm:main} from the prime number theorem \cite[pp.~6--16]{landau-thesis}; Landau published the converse implication in 1911 \cite[pp.~973--979]{landau-equivalence}. See also the elegant textbook treatment by Shapiro \cite[\S10.3, pp.~408--414]{shapiro}.) Here ``equivalent'' means that if either statement is assumed, there is a quick proof of the other, significantly shorter than a proof of either statement from scratch. The adverb ``elementarily'' draws a contrast with the more typical complex function-theoretic methods employed in prime number theory: The quick proofs in question involve only real variables, with Riemann's zeta function never (explicitly) mentioned.

It was quite a sensation when, in 1948, Erd\H{o}s and Selberg
announced that the prime number theorem itself could be established by
elementary means \cite{erdos,selberg}. For historical context, see Spencer and Graham’s article \cite{spencer-graham}, which reproduces notes by Ernst Straus, and Goldfeld’s account \cite{goldfeld}, which also draws on letters from files Selberg made available to him (cf.~the interview \cite{BS08}).
Both Erd\H{o}s's and Selberg's arguments use Selberg's ``symmetry
formula''; see \cite[eq.~(10.4.10), p.~419]{shapiro}.
In 1984, Daboussi \cite{daboussi} gave a very different elementary
proof of the PNT, using smooth numbers. Further elementary proofs
were later given by Hildebrand \cite{hildebrand}, via a large sieve
inequality, and --- quite recently --- Richter \cite{richter}, through a general result
on the distribution of $\Omega(n)$, the number of prime factors
of $n$ counted with multiplicity.

In this note, we present a simplification of Daboussi's argument for Theorem~\ref{thm:main}. Let $P^{+}(n)$ denote the largest prime factor of $n$, adopting the convention that $P^{+}(1):=1$. We say $n$ is $y$-smooth if $P^{+}(n)\le y$. For $x\ge1$ and $y\ge 2$, put
\[ \Psi(x,y):=\sum_{\substack{n\le x \\ P^{+}(n)\le y}} 1, \qquad M(x,y):=\sum_{\substack{n\le x \\ P^{+}(n)\le y}} \mu(n).\]
Set \[ a:=\limsup_{x\to\infty}\frac{|M(x)|}{x},\] so that the PNT is equivalent to proving that $a=0$. 

A central component of Daboussi's proof is a comparison between $|M(x,y)|$ and $a\Psi(x,y)$, based on slight variants of the ``logarithmic'' relations
\begin{equation}\label{eq:mobius-log0}
|M(x,y)|\log x
\le \sum_{\substack{p\le y\\p^k\le x}}\log p\,|M(x/p^k,y)|
 +\int_1^x |M(t,y)|\dt,
\end{equation}
\begin{equation}\label{eq:psi-log0}
\Psi(x,y)\log x
=\sum_{\substack{p\le y\\p^k\le x}}\log p\,\Psi(x/p^k,y)
 +\int_1^x\Psi(t,y)\dt.
\end{equation}
(We prove both relations in \S\ref{sec:excess} below.)
Daboussi extracts what he needs from \eqref{eq:mobius-log0} and
\eqref{eq:psi-log0} by integrating against the same carefully
constructed kernel function; see \cite[Lemma~2, p.~162, and
\S III.3 and \S IV, pp.~163--164]{daboussi}. 
We avoid any need to introduce a special kernel function by working
directly with the inequality for
$\bigl(|M(x,y)|-a\Psi(x,y)\bigr)\log x$ that results from subtracting
$a$ copies of \eqref{eq:psi-log0} from \eqref{eq:mobius-log0}.

In their lovely book \cite{tmf}, Tenenbaum and Mend\`es France describe a variant of Daboussi's argument that, like ours, avoids Daboussi's kernel function. The tradeoff is the need to develop certain ground-floor estimates from the theory of smooth numbers, including that $\Psi(x^u,x)/x^u \approx \rho(u)$ (Dickman's $\rho$-function) for bounded $u$ (with a key step in their proof being that $|M(x^u,x)|/x^u \lessapprox a\rho(u)$ in the same range). Our arrangement of the argument requires no nontrivial estimates at all on $\Psi(x,y)$.

% Daboussi's original argument avoids these estimates for
% $\Psi(x,y)$. Instead, he constructs a positive, decreasing,
% smooth function $k(s)$ satisfying \cite[Lemma~2]{daboussi}
% \[
% s k(s)-\int_s^{s+1}k(v)\,\mathrm{d}v=1\qquad(s>0).
% \]
% The construction is elementary,
% although the choice of kernel is not immediately apparent.
% Our argument dispenses with both the smooth-number estimates
% and the kernel. We obtain the required bound for the excess
% directly from the logarithmic identities for $M(x,y)$ and
% $\Psi(x,y)$, using Mertens' estimates.

Of course, we do need \emph{something}. We will make free use of the following two estimates of Mertens, valid for $x\ge2$:
\[
\sum_{p\le x}\frac{\log p}{p}=\log x+O(1),\qquad
\prod_{p\le x}\left(1-\frac1p\right)^{-1}\ll\log x.
\]
Here and throughout, the variable $p$ is reserved for primes. For (elementary)
proofs, see Shapiro \cite[Chapter~9]{shapiro}.

\subsection*{Use of LLMs} This is a work of exposition. It is offered as a (hopefully!) human-readable ``digestion'' of output produced by ChatGPT 6 Astra, as prompted by the author. P.\,P. takes no credit for the new mathematics but stands with, by, and beside the final manuscript.

Lean formalizations of Theorem~\ref{thm:main} and the smooth-number argument presented here, together with verification records and instructions for reproducing the checks, are available at \url{https://github.com/twonth/daboussi-pnt-lean}. The formalized results use only Lean's standard axioms of propositional extensionality (\texttt{propext}), classical choice (\texttt{Classical.choice}), and quotient soundness (\texttt{Quot.sound}).

P.\,P.~thanks OpenAI and Anthropic for providing complimentary access through the ChatGPT for Academic Researchers program and the Claude Team plan for scientists, respectively.

\subsection*{Other acknowledgements} P.\,P. presented this argument in a Fall 2026 graduate course offered at the University of Georgia. He thanks the class, and the university, for allowing him this opportunity to workshop the proof.

\section{A Fundamental Lemma and the deduction of Theorem \ref{thm:main}}\label{sec:fundamental}
The M\"{o}bius function is important in number theory precisely because it provides an elegant, number-theoretically friendly answer to the question of ``who is number $1$'': 
\[ \mathbf1_{n=1} = \sum_{d\mid n}\mu(d).\] Letting $x\ge 1$ and summing this identity over $n\le x$,
\[
1=\sum_{d\le x}\mu(d)\left\lfloor\frac xd\right\rfloor
 =x\sum_{d\le x}\frac{\mu(d)}d
  -\sum_{d\le x}\mu(d)\left\{\frac xd\right\}.
\]
Consequently,
\[
\left|\sum_{d\le x}\frac{\mu(d)}d\right|
<\frac{1+\lfloor x\rfloor}{x}\le2\qquad(x\ge1).
\]
Expressing the same sum by partial summation, we have
\[
\sum_{n\le x}\frac{\mu(n)}n
=\frac{M(x)}x+\int_1^x\frac{M(t)}t\dt,
\]
and hence
\begin{equation}\label{eq:signed}
\left|\int_1^x\frac{M(t)}t\dt\right|<3\qquad(x\ge1).
\end{equation}

Our proof of Theorem \ref{thm:main} combines \eqref{eq:signed} with the following estimate, whose
proof occupies \S\S\ref{sec:smooth}--\ref{sec:excess}.
\begin{lemma}[Fundamental Lemma]\label{lem:fundamental}
As $x\to\infty$,
\[
\int_1^x\pos{a-\frac{|M(t)|}{t}}\dt=o(\log x).
\]
{\rm (}Here $v^+=\max\{v,0\}$ denotes the positive part of $v$.{\rm )}
\end{lemma}

Roughly speaking, Lemma \ref{lem:fundamental} says that $|M(t)|/t$ spends most of its time close to its limit superior $a$, with time measured logarithmically.

\begin{proof}[Proof of Theorem~\ref{thm:main}, assuming the Fundamental Lemma]
Suppose for a contradiction that $a>0$, and fix $B:=\exp(12/a)$.
We first use \eqref{eq:signed} to show that every interval $[X,BX]$,
with $X\ge1$, contains a $u$ for which
\begin{equation}\label{eq:small}
|M(u)|\le\frac a2u.
\end{equation}
Indeed, if there is no such $u$, then $M(t)$ is nonzero
throughout $[X,BX]$. Since $M$ takes integer values and jumps by
at most $1$, it has constant sign there. This gives
\[
\left|\int_X^{BX}\frac{M(t)}t\dt\right|
=\int_X^{BX}\frac{|M(t)|}t\dt
\ge\frac a2\log B=6.
\]
But \eqref{eq:signed}, applied at $BX$ and at $X$, bounds the
left-hand side strictly below $6$. 

Now let $X\ge X_0:=8/a$, and choose $u\in[X,BX]$ satisfying
\eqref{eq:small}. To obtain a lower bound for the integral in the
Fundamental Lemma, we extend this estimate to nearby values of $t$.
For $u\le t\le u(1+a/8)$, the bound $|\mu(n)|\le1$ gives
\[
\begin{aligned}
|M(t)|
&\le |M(u)|+|M(t)-M(u)|\\
&\le\frac a2u+(t-u)+1\\
&\le\frac a2u+\frac a8u+\frac a8u
\le\frac{3a}{4}t,
\end{aligned}
\]
where we used $1\le au/8$. The bounds $u\le BX$ and $a\le1$ also
give $u(1+a/8)\le2BX$, so
\[
\int_X^{2BX}\pos{a-\frac{|M(t)|}t}\dt
\ge\int_u^{u(1+a/8)}\frac a4\dt
=\frac a4\log(1+a/8)>0.
\]
We apply this bound with $X=(2B)^{j} X_0$, for $j=0,\ldots, N-1$, noting that the corresponding intervals $(X,2BX]$ are disjoint and contained in $(1,(2B)^{N} X_0]$. Adding these estimates gives
\[
\int_1^{(2B)^NX_0}\pos{a-\frac{|M(t)|}t}\dt
\ge N\frac a4\log(1+a/8).
\]
But the Fundamental Lemma implies that our left-hand side is $o(\log((2B)^NX_0))=o(N)$ as $N\to\infty$. This contradiction forces $a=0$.
\end{proof}

\section{Splitting off bounded prime factors}\label{sec:smooth}
Recall from the introduction that $\Psi(x,y)=\sum_{n\le x,\,P^{+}(n)\le y} 1$ while $M(x,y)=\sum_{n\le x,\,P^{+}(n)\le y} \mu(n)$. Fix $y\ge 2$, and put
\[
P=\prod_{p\le y}p.
\]
We have $|M(t,y)|\le\Psi(t,y)$ and, since every positive integer at
most $y$ is $y$-smooth,
\[
M(t,y)=M(t),\qquad \Psi(t,y)=\lfloor t\rfloor\qquad(1\le t\le y).
\]
At the other end of the range, only the squarefree $y$-smooth
integers contribute to $M(t,y)$, and these are precisely the
divisors of $P$. Thus, for $t\ge P$,
\begin{equation}\label{eq:support}
M(t,y)=\sum_{d\mid P}\mu(d)=0.
\end{equation}

Following Daboussi, we compare $|M(x,y)|$ with $a\Psi(x,y)$. The following estimate gets us started.

\begin{proposition}\label{prop:comparison}
For every $y\ge2$,
\begin{equation}\label{eq:comparison}
\int_1^\infty
\left(\frac{|M(t,y)|}t-a\frac{\Psi(t,y)}t\right)\dt\ge0.
\end{equation}
\end{proposition}
\begin{proof} Each positive integer admits a unique decomposition in the form $n=bc$, where
$P^+(b)\le y$ and $(c,P)=1$. Since $(b,c)=1$, the multiplicativity
of $\mu$ gives
\[
M(x)=\sum_{\substack{c\le x\\(c,P)=1}}\mu(c)M(x/c,y).
\]
The terms with $c\le x/P$ vanish by \eqref{eq:support}. Taking
absolute values in the remaining terms and grouping according to the value of
$j:=\lfloor x/c\rfloor<P$, we obtain
\[
|M(x)|\le
\sum_{j=1}^{P-1}|M(j,y)|
\sum_{\substack{x/(j+1)<c\le x/j\\(c,P)=1}}1.
\]
View $y$ as fixed. Since every block of $P$ consecutive integers $c$ contains precisely $\varphi(P)$ that satisfy $(c,P)=1$, 
\[
\sum_{\substack{x/(j+1)<c\le x/j\\(c,P)=1}}1
=x\frac{\varphi(P)}P\left(\frac1j-\frac1{j+1}\right)+O_y(1).
\]
Substituting this into our bound for $|M(x)|$, dividing by $x$,
and letting $x\to\infty$, we find that
\[
\begin{aligned}
a
&\le\frac{\varphi(P)}P
 \sum_{j=1}^{P-1}|M(j,y)|\left(\frac1j-\frac1{j+1}\right)\\
&=\frac{\varphi(P)}P
 \int_1^\infty\frac{|M(t,y)|}t\dt.
\end{aligned}
\]
(For the final equality, we use that $M(t,y)=M(j,y)$ on $[j,j+1)$
and $M(t,y)=0$ for $t\ge P$.) Rearranging, we have that
\begin{equation}\label{eq:smooth-lower}
\int_1^\infty\frac{|M(t,y)|}t\dt
\ge a\frac P{\varphi(P)}
=a\prod_{p\le y}(1-1/p)^{-1}.
\end{equation}

The product on the right also occurs in the integral of
$\Psi(t,y)/t$. In fact,
\begin{equation}\label{eq:positive-mass}
\int_1^\infty\frac{\Psi(t,y)}t\dt
=\sum_{P^+(n)\le y}\frac1n
=\prod_{p\le y}(1-1/p)^{-1}.
\end{equation}
To see this, observe that by partial summation,
\begin{equation}\label{eq:smoothharmonic}
\sum_{\substack{n\le X\\P^+(n)\le y}}\frac1n
=\frac{\Psi(X,y)}X+\int_1^X\frac{\Psi(t,y)}t\dt.
\end{equation}
For every positive integer $n\le X$, and every prime $p$, the exponent of $p$ in $n$ is at most $\log{X}/\log{2}$. Hence (keeping in mind that $y$ is fixed), 
\[
\Psi(X,y)\le
\left(1+\frac{\log X}{\log2}\right)^{\pi(y)}=o(X).
\]
Sending $X$ to infinity in \eqref{eq:smoothharmonic} yields the first equality in 
\eqref{eq:positive-mass}. The second follows by expanding the Euler product. Finally, the proposition itself follows by comparing \eqref{eq:smooth-lower} and \eqref{eq:positive-mass}.\end{proof}

To apply Proposition~\ref{prop:comparison} to the Fundamental Lemma,
we separate its integrand into positive and negative parts. Define ``excess'' and ``deficit'' functions by
\[
E_y(t)=\pos{\frac{|M(t,y)|}t-a\frac{\Psi(t,y)}t},
\qquad
D_y(t)=\pos{a\frac{\Psi(t,y)}t-\frac{|M(t,y)|}t}.
\]
% Both are integrable with respect to $\mathrm{d}t/t$ by
% \eqref{eq:positive-mass}, since
% \[
% E_y(t)\le\frac{\Psi(t,y)}t,\qquad
% D_y(t)\le a\frac{\Psi(t,y)}t.
% \]
The integrand in \eqref{eq:comparison} is $E_y(t)-D_y(t)$, so that Proposition~\ref{prop:comparison} can be restated as the assertion that
\begin{equation}\label{eq:balance}
\int_1^\infty D_y(t)\dt\le\int_1^\infty E_y(t)\dt.
\end{equation}
That is, the total excess of $|M(t,y)|/t$ over $a\Psi(t,y)/t$ is at least the total deficit, both measured with respect to $\mathrm{d}t/t$.

When $1\le t\le y$, the deficit $D_y(t)$ is closely related to the integrand in the Fundamental Lemma:
\[
\begin{aligned}
D_y(t)
&=\pos{a\frac{\lfloor t\rfloor}t-\frac{|M(t)|}t}\\
&=\pos{a-\frac{|M(t)|}t-a\frac{\{t\}}t}
\ge\pos{a-\frac{|M(t)|}t}-\frac at.
\end{aligned}
\]
(We use here that $(v-h)^+\ge v^+-h$ for $h\ge0$.) Integrating with respect to $\mathrm{d}t/t$, and noting that $\int_1^{y} a/t\,\mathrm{d}t/t \le a$, we obtain
\begin{equation*}
\int_1^y\pos{a-\frac{|M(t)|}t}\dt
\le\int_1^yD_y(t)\dt+a
\le\int_1^\infty E_y(t)\dt+a,
\end{equation*}
using \eqref{eq:balance} in the last step.
We have therefore reduced the Fundamental Lemma to proving that,
as $y\to\infty$,
\begin{equation}\label{eq:excess-goal}
\int_1^\infty E_y(t)\dt=o(\log y).
\end{equation}

\section{Making less of the excess}\label{sec:excess}

\begin{proof}[Proof of Lemma~\ref{lem:fundamental}]
To prove \eqref{eq:excess-goal}, we start with the trivial bounds
\[
0\le E_y(t)\le\frac{|M(t,y)|}t\le\frac{\Psi(t,y)}t\le1.
\]
Together with \eqref{eq:positive-mass} and Mertens' estimate, this
gives
\begin{equation}\label{eq:excess-bigO}
\int_1^\infty E_y(t)\dt
\le\prod_{p\le y}(1-1/p)^{-1}\ll\log y.
\end{equation}
So the remaining fight is over the case of a single letter: We must improve $O(\log y)$ to $o(\log y)$.

We split the integral at $y$. The required estimate for $\int_1^{y} E_y(t)\,\mathrm{d}t/t$ follows from the definition of
$a$. Given $\eps>0$, choose $T\ge1$ so that $|M(t)|/t\le a+\eps$
for $t\ge T$. Then, for $T\le t\le y$,
\[
E_y(t)
=\pos{\frac{|M(t)|}t-a+a\frac{\{t\}}t}
\le\eps+\frac at.
\]
For $y\ge T$, integrating this estimate, using that $E_y(t)\le1$ below $T$, we get that
\[
\int_1^yE_y(t)\dt
\le\log T+\eps\log(y/T)+a.
\]
Hence,
\[  \limsup_{y\to\infty} \frac{1}{\log{y}}\int_1^yE_y(t)\dt \le \eps.\]
As $\eps > 0$ is arbitrary,
\begin{equation}\label{eq:below}
\int_1^yE_y(t)\dt=o(\log y).
\end{equation}

To handle the remaining range $(y,\infty)$, we need the ``logarithmic'' relations quoted in the introduction. We begin their derivation with some simple consequences of M\"{o}bius inversion. Writing $\Lambda$ for the von Mangoldt function and starting from
\begin{equation}\label{eq:chebyshevlog} \log{n} = \sum_{p^k \mid n} \log{p} = \sum_{d \mid n} \Lambda(d),  \end{equation}
we find that
\begin{align*} \Lambda(n) &= \sum_{d\mid n} \mu(d)\log{(n/d)} 
\\&= \log{n} \sum_{d\mid n} \mu(d) - \sum_{d \mid n}\mu(d)\log{d}\\
&=- \sum_{d \mid n}\mu(d)\log{d}. \end{align*}
Inverting this last formula for $\Lambda(n)$ gives
\begin{equation}\label{eq:muchebyshev}
-\mu(n)\log n=\sum_{p^k\mid n}\log p\,\mu(n/p^k).
\end{equation}
(Here and below, sums over $p^k$ run over primes $p$ and positive
integers $k$.) 

Summing \eqref{eq:muchebyshev} over $y$-smooth integers $n\le x$, and reversing the order of summation, 
\[
\begin{aligned}
-\sum_{\substack{n\le x\\P^+(n)\le y}}\mu(n)\log n
&=\sum_{\substack{p\le y\\p^k\le x}}\log p
  \sum_{\substack{m\le x/p^k\\P^+(m)\le y}}\mu(m)\\
&=\sum_{\substack{p\le y\\p^k\le x}}\log p\,M(x/p^k,y).
\end{aligned}
\]
On the other hand, partial summation gives
\[
\sum_{\substack{n\le x\\P^+(n)\le y}}\mu(n)\log n
=M(x,y)\log x-\int_1^xM(t,y)\dt.
\]
Substituting this into the immediately preceding display and rearranging,
\begin{equation*}
M(x,y)\log x
=-\sum_{\substack{p\le y\\p^k\le x}}\log p\,M(x/p^k,y)
 +\int_1^xM(t,y)\dt.
\end{equation*}
Taking absolute values and appealing to the triangle inequality,
\begin{equation}\label{eq:mobius-log2}
|M(x,y)|\log x
\le \sum_{\substack{p\le y\\p^k\le x}}\log p\,|M(x/p^k,y)|
 +\int_1^x |M(t,y)|\dt.
\end{equation}
This is nothing other than the relation \eqref{eq:mobius-log0}. Carrying out the same calculation with \eqref{eq:chebyshevlog} replacing \eqref{eq:muchebyshev} gives
\begin{equation}\label{eq:psi-log}
\Psi(x,y)\log x
=\sum_{\substack{p\le y\\p^k\le x}}\log p\,\Psi(x/p^k,y)
 +\int_1^x\Psi(t,y)\dt,
\end{equation}
which appeared earlier as \eqref{eq:psi-log0}.

Subtracting $a$ times \eqref{eq:psi-log} from \eqref{eq:mobius-log2}, we obtain
\begin{align*}
\bigl(|M(x,y)|-a\Psi(x,y)\bigr)\log x
&\le\sum_{\substack{p\le y\\p^k\le x}}\log p
 \bigl(|M(x/p^k,y)|-a\Psi(x/p^k,y)\bigr)\\
&\qquad+\int_1^x\bigl(|M(t,y)|-a\Psi(t,y)\bigr)\dt.
\end{align*}
Replacing each difference on the right by its positive part only
increases the right-hand side, and so preserves the inequality. Since the resulting expression is nonnegative, the inequality is still preserved if we then replace the left-hand side with \emph{its} positive part. Noting that $(|M(t,y)|-a\Psi(t,y))^+=tE_y(t)$, we deduce after dividing by $x$ that 
\begin{equation}\label{eq:excess-inequality}
E_y(x)\log x
\le\sum_{\substack{p\le y\\p^k\le x}}
 \frac{\log p}{p^k}E_y(x/p^k)
 +\int_1^x\frac tx E_y(t)\dt.
\end{equation}

We integrate \eqref{eq:excess-inequality} with respect to $\mathrm{d}x/x$, 
over $[1,\infty)$. For each prime power $p^k$, substituting
$x=p^kt$ gives
\[
\int_{p^k}^\infty\frac{\log p}{p^k}E_y(x/p^k)\dx
=\frac{\log p}{p^k}\int_1^\infty E_y(t)\dt.
\]
For the remaining term, we reverse the order of integration:
\[
\begin{aligned}
\int_1^\infty\left(\int_1^x\frac tx E_y(t)\dt\right)\dx
&=\int_1^\infty E_y(t)
 \left(\int_t^\infty\frac tx\dx\right)\dt\\
&=\int_1^\infty E_y(t)\dt.
\end{aligned}
\]
(Here all manipulations of sums and integrals are justified by nonnegativity.) Combining the two contributions, 
\[
\int_1^\infty E_y(t)\log t\dt
\le\left(1+\sum_{p\le y}\sum_{k\ge1}\frac{\log p}{p^k}\right)
\int_1^\infty E_y(t)\dt.
\]
The coefficient on the right is estimated using Mertens' theorem:
\[
\begin{aligned}
\sum_{p\le y}\sum_{k\ge1}\frac{\log p}{p^k}
&=\sum_{p\le y}\frac{\log p}{p}
  +\sum_{p\le y}\frac{\log p}{p(p-1)}
=\log y+O(1).
\end{aligned}
\]
Hence, for some absolute
constant $C>0$,
\[
\int_1^\infty E_y(t)\log t\dt
\le(\log y+C)\int_1^\infty E_y(t)\dt.
\]
Rearranging and applying \eqref{eq:excess-bigO}, we find that
\begin{equation}\label{eq:centered}
\int_1^\infty E_y(t)\log(t/y)\dt
\le C\int_1^\infty E_y(t)\dt=O(\log y).
\end{equation}
We are almost home! We already had an $O(\log y)$ bound without
the factor $\log(t/y)$, but this factor is sensitive to \emph{where}
the excess occurs.

Recall we have already shown that $\int_1^{y} E_y(t)\, \mathrm{d}t/t = o(\log{y})$. Slightly larger values of $t$, say $y\le t \le y^{1+\eps}$ with $\eps > 0$ fixed, pose no trouble. Since $E_y(t)\le1$,
\[
\int_y^{y^{1+\eps}}E_y(t)\dt
\le\int_y^{y^{1+\eps}}\dt
=\eps\log y,
\]
which will be acceptable to us upon taking $\eps$ small.

It remains to consider $t$ that are substantially larger than $y$, meaning $t\ge y^{1+\eps}$. Moving the contribution from
$[1,y]$ in \eqref{eq:centered} to the right and using
\eqref{eq:below}, we obtain
\begin{align}
\int_y^\infty E_y(t)\log(t/y)\dt
&\le\int_1^y E_y(t)\log(y/t)\dt+O(\log y)\notag\\
&\le(\log y)\int_1^y E_y(t)\dt+O(\log y)\notag\\
&=o((\log y)^2).\label{eq:weighted-tail}
\end{align}
On the other hand, $\log(t/y)\ge\eps\log y$ when
$t\ge y^{1+\eps}$, so
\[
\int_y^\infty E_y(t)\log(t/y)\dt
\ge\int_{y^{1+\eps}}^\infty E_y(t)\log(t/y)\dt
\ge\eps\log y\int_{y^{1+\eps}}^\infty E_y(t)\dt.
\]
Comparing with \eqref{eq:weighted-tail}, we find that
\[
\int_{y^{1+\eps}}^\infty E_y(t)\dt=o(\log y).
\]
Combining the estimates for the three ranges gives
$\int_1^\infty E_y(t)\,\mathrm{d}t/t 
\le\eps\log y+o(\log y)$. Dividing by $\log y$ and letting $y\to\infty$, we obtain
\[
\limsup_{y\to\infty}\frac1{\log y}
\int_1^\infty E_y(t)\dt\le\eps.
\]
But $\eps>0$ is arbitrary, and so $\int_1^\infty E_y(t) \,\mathrm{d}t/t=o(\log y)$, as desired. This proves \eqref{eq:excess-goal}, and hence the Fundamental
Lemma. As shown in \S\ref{sec:fundamental}, it follows that $M(x)=o(x)$.
\end{proof}

\begin{thebibliography}{99}

\bibitem{BS08}
N.~A. Baas and C.~F. Skau, \emph{The lord of the numbers, Atle Selberg. On his life and mathematics}, Bull. Amer. Math. Soc. (N.S.) \textbf{45} (2008), no.~4, 617--649.

\bibitem{daboussi}
H. Daboussi, \emph{Sur le th\'eor\`eme des nombres premiers},
C. R. Acad. Sci. Paris S\'er. I Math. \textbf{298} (1984),
no.~8, 161--164.

\bibitem{erdos}
P. Erd\H{o}s, \emph{On a new method in elementary number theory
which leads to an elementary proof of the prime number theorem},
Proc. Nat. Acad. Sci. U.S.A. \textbf{35} (1949), 374--384.

\bibitem{goldfeld}
D. Goldfeld, \emph{The elementary proof of the prime number theorem:
an historical perspective}, in \emph{Number Theory} (New York, 2003),
D. Chudnovsky, G. Chudnovsky, and M. Nathanson (eds.),
Springer, New York, 2004, 179--192.

\bibitem{hildebrand}
A. Hildebrand, \emph{The prime number theorem via the large sieve},
Mathematika \textbf{33} (1986), 23--30.

\bibitem{landau-thesis}
E. Landau, \emph{Neuer Beweis der Gleichung
$\sum_{k=1}^{\infty}\mu(k)/k=0$}, Inauguraldissertation,
Berlin, 1899.

\bibitem{landau-equivalence}
E. Landau, \emph{\"Uber die \"Aquivalenz zweier Haupts\"atze
der analytischen Zahlentheorie}, Sitzungsber. Kaiserl. Akad.
Wiss. Wien, Math.-Naturwiss. Kl., Abt.~IIa \textbf{120} (1911),
973--988.

\bibitem{mangoldt}
H. von Mangoldt, \emph{Beweis der Gleichung
$\sum_{k=1}^{\infty}\mu(k)/k=0$}, Sitzungsber. K\"onigl.
Preu\ss. Akad. Wiss. Berlin (1897), 835--852.

\bibitem{richter}
F. K. Richter, \emph{A new elementary proof of the Prime Number
Theorem}, Bull. London Math. Soc. \textbf{53} (2021), 1365--1375.

\bibitem{selberg}
A. Selberg, \emph{An elementary proof of the prime-number theorem},
Ann. of Math. (2) \textbf{50} (1949), 305--313.

\bibitem{shapiro}
H. N. Shapiro, \emph{Introduction to the Theory of Numbers},
John Wiley \& Sons, New York, 1983; reprinted by Dover
Publications, Mineola, NY, 2008.

\bibitem{spencer-graham}
J. Spencer and R. Graham, \emph{The elementary proof of the prime
number theorem}, Math. Intelligencer \textbf{31} (2009),
no.~3, 18--23.

\bibitem{tmf}
G. Tenenbaum and M. Mend\`es France, \emph{The Prime Numbers and
Their Distribution}, translated by P. G. Spain, Student
Mathematical Library, vol.~6, American Mathematical Society,
Providence, RI, 2000.
\end{thebibliography}
\end{document}
